\subsection{A bounded list of dyadic supersingular comparison levels}

\begin{lemma}[Dyadic graph and logarithm estimates]
\label{lem:B-RD45}
\status{Candidate; complete independent audit requested.}
Let $K$ have algebraically closed residue field of characteristic two.
Suppose $M,N$ are height-two potentially supersingular groups with
different prime-to-two special-inertia characters and
$M[2^n]\simeq N[2^n]$. Let $L/K$ be a common finite Galois
good-reduction extension, put $R=\mathcal O_L$ and $a=v_L(2)$,
and write $M_0,N_0$ for the special good groups.
Define
\[
 q_0=1,\qquad q_{r+1}=\min\{a+q_r,1+2q_r,4q_r\}.
 \tag{B-RD45.1}
\]
Then $q_{n-1}\le a$.
More precisely, the special graph at every level $2^r$, $r\le n$,
is $\ker F_{M_0}^r\times\ker F_{N_0}^r$; if $q_u>a$, then
$u\ge n$.
\end{lemma}
\begin{proof}
We first supply the dyadic replacement for the tame idempotent
in Proposition~\ref{prop:B-RD33}.
By Lemma~\ref{lem:B-RD44-conjugacy}, exactly one of the two
faithful finite inertia orders is divisible by three.
Choose an element $\tau$ of order three in the joint finite image.
On one special group it is the identity; on the other it has
minimal polynomial $T^2+T+1$ in the quaternion endomorphism algebra.
Consequently
\[
                       e_0=(1+\tau+\tau^2)/3
 \tag{B-RD45.2}
\]
is the integral projection onto the first factor. Division by
three is valid on every two-primary group scheme.

By Proposition~\ref{prop:B-RD44-realization} the generic groups
are exact elliptic BT groups and have perfect alternating
pairings into $\mu_{2^\infty}$. Alternatively their finite
determinants are one and the determinant argument in B-RD33
gives the same pairings. Tate's integral full faithfulness
extends each pairing and its inverse to the good models.
At level $2^r$, rescale one factor by the unit multiplier of
the asserted rank-two isomorphism. Its graph is then maximal
isotropic. Its finite flat closure $H_r$ remains maximal
isotropic after reduction: restricted pairings extend by
flatness, finite-flat duality commutes with reduction, and
the closed inclusion into its orthogonal has equal free rank.
These are the exact arguments of B-RD33, with no division by two.

The idempotent (B-RD45.2) preserves $(H_r)_k$ and splits it
as a product of subgroups of $M_0,N_0$.
In a one-dimensional connected formal group the unique subgroup
of rank $2^b$ is $\ker F^b$, since its defining ideal in
$k[[T]]$ must be $(T^{2^b})$.
The two ranks have exponents summing to $2r$; self-orthogonality
in each factor makes each exponent $r$.
Thus
\[
                 (H_r)_k=\ker F_{M_0}^r\times\ker F_{N_0}^r.
 \tag{B-RD45.3}
\]

We give the logarithm argument in the required prime-two form.
Let $G/R$ be any finite-height formal group whose special
two-series has no monomial of total degree less than four.
This condition holds for a product of one-dimensional
height-two groups and is invariant under integral changes
of formal coordinates. Its linear term is $2X$; coefficients
of total degrees two and three are in $\pi R$, and all higher
coefficients are integral. Therefore $[2]_G$ maps coefficient
valuation at least $q$ to valuation at least
$\min(a+q,1+2q,4q)$.
For any augmentation vector $h\in tR[[t]]^{\dim G}$,
the coefficients of $[2]^u_G(\pi h)$ have valuations at
least $q_u$.

Write $c_\alpha$ for a coefficient of total degree $m$ of
the normalized logarithm. Its differential is the integral
matrix of invariant differentials, so
\[
 v(c_\alpha)\ge
 -a\min_{j:\alpha_j>0}v_2(\alpha_j)
 \ge-a\lfloor\log_2m\rfloor.
\]
Since $\lfloor\log_2m\rfloor\le m-1$, substitution of a
vector of coefficient valuation $q>a$ gives logarithm
coefficient valuation at least $q$. All substitutions are
coefficientwise finite because the vectors have zero constant
term. Consequently
\[
        2^u\log_G(\pi h)\in\pi^{q_u}R[[t]]^{\dim G}
                  \quad\text{if }q_u>a.
 \tag{B-RD45.4}
\]

The Cartier-module proof of Lemma~\ref{lem:B-RD34-kernel}
now applies literally with prime two, with
\[
 \mathcal R=R[[\Delta]][1/2],\qquad J=\pi^aR=2R.
\]
To specify the prime-two point completely, the published original
explicitly imposes no $p>2$ condition
\cite[p.~663]{Bondarko2006Generic}.
Its Propositions~9--11 and Lemma~3, pp.~674--675, give
$J[[\Delta]]^d\subset D_G$, the identity
$\mathcal C_G=\Omega\pi D_{G_\pi}$, and the bounded-projection
embedding. The integral Witt-coefficient action of $\Omega$
preserves coefficient valuation ideals, so (B-RD45.4) gives
$2^u\mathcal C_G\subset J\{\{\Delta\}\}^d$.
Equation~(8), p.~676, gives
\[
 \mathbf f^uAz\equiv2^uAz\Delta^{-u}
                              \pmod{2R((\Delta))^b}.
\]
Because $2R=J$, an integral matrix $A$ and a bounded projection
$Az$ therefore satisfy $\mathbf f^uAz\in J[[\Delta]]^b$.
The saturation argument of Proposition~16, p.~681, and
Theorem~2(8), p.~668, then show that $F^u$ kills the special
kernel of any formal homomorphism having no nonzero finite
flat subgroup in its kernel, finite special kernel, and target
$G$ of this type. This is exactly the proof written in
B-RD34, including its distinction between absence of a flat
kernel and a finite-level closed immersion; no additional
prime-two full-faithfulness claim is inserted.

Finally set $P=\mathcal M\times\mathcal N$, $Q=P/H_n$, and
consider
\[
        P\longrightarrow Q\times\mathcal M,
        \qquad z\longmapsto(qz,\operatorname{pr}_{\mathcal M}z).
\]
It has no nonzero finite flat kernel, since generically the
graph of an isomorphism meets $0\times N$ trivially.
By (B-RD45.3) its special kernel is
$0\times\ker F_{N_0}^n$.
The target special fibre is a product of three height-two
formal groups: the quotient special group is the product of
the two Frobenius twists, because finite-flat quotient formation
commutes with base change. If $q_u>a$, the preceding estimate
would make $F^u$ annihilate this special kernel.
For $u<n$ its image has rank $2^{n-u}>1$, a contradiction.
Taking $u=n-1$ proves the asserted inequality.
\end{proof}

\begin{theorem}[B-RD45: a finite set of possible dyadic least indices]
\label{thm:B-RD45}
\status{Candidate; complete independent audit requested.}
For the manuscript question with $(p,e,h)=(2,1,2)$,
\[
                         3\le\nu_{\min}(2,1,2)\le7.
 \tag{B-RD45.5}
\]
More precisely a failure at the indicated torsion level must
compare two potentially supersingular elliptic BT groups after
completed unramified base change, with joint faithful finite
inertia degree in the corresponding row:
\[
\begin{array}{c|c|c}
\text{torsion level}&\text{manuscript index}&\text{remaining joint degrees}\\
\hline
16&3&24,48,96,192\\
32&4&48,96,192\\
64&5&96,192\\
128&6&192\\
256&7&\text{none}.
\end{array}
 \tag{B-RD45.6}
\]
The table is a necessary-condition table. It asserts no
counterexample in any remaining row, and does not decide
which of the five possible least indices is attained.
\end{theorem}
\begin{proof}
The lower bound is Proposition~\ref{prop:B-RD43}.
Proposition~\ref{prop:B-RD44} excludes every mixed case from
level sixteen onward and gives joint degrees $m=3\cdot2^r$,
$0\le r\le6$. Since $e_K=1$ and the residue field is algebraically
closed, a faithful joint good-reduction extension has
$a=v_L(2)=m$.
For $q\ge1$ the recurrence in (B-RD45.1) is simply
$q_{i+1}=\min(a+q_i,1+2q_i)$.
Starting from one, its first crossing of $a$ occurs at the
first integer $u$ with $2^{u+1}-1>a$.
For $a=3\cdot2^r$ this is $u=r+1$, since
$2^{r+1}-1<a<2^{r+2}-1$.
Lemma~\ref{lem:B-RD45} requires $n\le r+1$ at level $2^n$.
For $n=4,5,6,7,8$ this gives exactly the displayed rows,
including impossibility at $n=8$ because $r\le6$.
The manuscript convention is $M(\nu)=M[2^{\nu+1}]$, which
gives the stated upper index seven.
\end{proof}

\begin{remark}[Scope of the arithmetic certificate]
The program \path{code/verify_p2_height_two_log_gate.py}
checks the cubic projector in the quaternion algebra, the
recurrence and first-crossing formula, the logarithm degree
inequalities, and every joint-degree entry of (B-RD45.6).
Its retained output records the exact scope and runtime.
The full Cartier saturation and graph arguments are proved
above, not inferred from the finite checks. In particular
no finite decision procedure for sixteen-congruences has
been executed.
\end{remark}
